%%%%%%%%%%%%%%%%%%%%%%%%%%%%%%%%%%%%%%%%%%%%%%%%%%%%%%%%%%%%%%%%%%%%%%%%%%%%%%
%  Unit II, second half: Search Engine Optimization (Chapters 6-11)
%%%%%%%%%%%%%%%%%%%%%%%%%%%%%%%%%%%%%%%%%%%%%%%%%%%%%%%%%%%%%%%%%%%%%%%%%%%%%%

\chapter{Search Engine Optimization Fundamentals}

\syllabusstrip{Search Engine Optimization: Search Engine Optimization Fundamentals.}

\begin{learningoutcomes}
\begin{itemize}
\item What SEO is, and its primary objective.
\item Why SEO matters --- the figures that justify the investment.
\item SEO, PPC and SEM distinguished.
\item The anatomy of a modern search engine results page.
\item The shift from keyword matching to intent, and E-E-A-T.
\item The six-step SEO process, and white-hat against black-hat technique.
\end{itemize}
\end{learningoutcomes}

\section{Definition and Objective}

\begin{definitionbox}[Search Engine Optimization]
\keytermas{Search engine optimization}{search engine optimization} (SEO) is the
set of techniques used to improve the visibility and ranking of a website in the
organic --- that is, unpaid --- results of search engines, thereby increasing the
quantity and quality of traffic to that site.
\end{definitionbox}

\begin{exambox}[Primary Objective --- 2 Marks]
\pyq{SEE, Jun/Jul 2025 --- 2 M}
The primary objective of SEO is to increase organic visibility for relevant
queries so that the right people --- those actively searching with intent ---
find the site, and to increase the site's usability\ix{usability} so that they are satisfied
when they arrive.

\smallskip
\textbf{Visibility without relevance is worthless.} Both halves of that sentence
are required for full marks.
\end{exambox}

\section{Why SEO Matters}

\begin{keybox}[The Case for Organic Search]
\begin{itemize}
\item Around \textbf{10 billion} searches happen every day on Google.
\item \textbf{53\%} of all website traffic is delivered by organic search\ix{organic search} alone.
\item Organic traffic has \textbf{infinite lifespan} --- it compounds, whereas
      paid traffic stops the moment the budget stops.
\item Search traffic carries the \textbf{highest buyer intent} of any channel,
      because the user has explicitly stated their need.
\end{itemize}
\end{keybox}

\noindent\textbf{Importance for online businesses:} lower customer acquisition
cost; credibility, because users trust organic results more than advertisements;
lead generation\ix{lead generation} around the clock; competitive defence; and measurable results.

\section{SEO, SEM and PPC --- The Visibility Matrix}

\begin{mmlongtable}{SEO, PPC and SEM distinguished}{tab:seosempcc}%
{@{}P{0.16\textwidth}P{0.46\textwidth}P{0.31\textwidth}@{}}
\rowcolor{ocre}\thd{Term} & \thd{Mechanism} & \thd{Cost model}\\
\midrule
\endfirsthead
\rowcolor{ocre}\thd{Term} & \thd{Mechanism} & \thd{Cost model}\\
\midrule
\endhead
\rlab{SEO} & Drives organic results and clicks & Free per click; costs time and
resources\\
\rowcolor{tablealt}
\rlab{PPC} \newline {\scriptsize pay per click} & Drives paid results;
advertisers bid on specific keywords & The advertiser is charged per click\\
\rlab{SEM} \newline {\scriptsize search engine marketing} & The umbrella term
combining both organic and paid search strategies & Both\\
\end{mmlongtable}

On a search engine results page, paid listings appear at the top and bottom, with
organic results between them.

\section{The Search Engine Results Page}

\begin{definitionbox}[SERP]
The \keytermas{SERP}{SERP} is the page a search engine returns in response to a
query. It is no longer a simple list of ten blue links.
\end{definitionbox}

\begin{figure}[htbp]
\centering
\begin{tikzpicture}[font=\sffamily\footnotesize,
    el/.style={minimum width=6.4cm,minimum height=0.62cm,align=left,
               text width=6.0cm,font=\sffamily\scriptsize,inner sep=4pt,
               anchor=north west}]
  \draw[rulegrey,line width=0.9pt,rounded corners=3pt] (0,0) rectangle (6.9,-8.5);
  \fill[tablehead,rounded corners=3pt] (0,0) rectangle (6.9,-0.55);
  \node[anchor=west,font=\sffamily\scriptsize,text=slatelight] at (0.18,-0.28)
    {search query};

  \node[el,fill=warnpale,draw=warnred] at (0.25,-0.72)
    {\textbf{\textcolor{warnred}{Paid results}} --- search and shopping ads};
  \node[el,fill=ocrepale,draw=ocre] at (0.25,-1.5)
    {\textbf{\textcolor{ocredark}{Featured snippet}} --- ``position zero''};
  \node[el,fill=deepbluepale,draw=deepblue] at (0.25,-2.28)
    {\textbf{\textcolor{deepblue}{Local pack}} --- three businesses, map, ratings};
  \node[el,fill=white,draw=rulegrey] at (0.25,-3.06)
    {\textbf{Organic result} --- title tag, URL, meta description};
  \node[el,fill=white,draw=rulegrey] at (0.25,-3.84)
    {\textbf{Organic result} --- with rich results: stars, FAQs, price};
  \node[el,fill=greenpale,draw=greenok] at (0.25,-4.62)
    {\textbf{\textcolor{greenok}{People Also Ask}} --- expandable questions};
  \node[el,fill=white,draw=rulegrey] at (0.25,-5.40)
    {\textbf{Image, video and news carousels}};
  \node[el,fill=white,draw=rulegrey] at (0.25,-6.18)
    {\textbf{Organic results} continue};
  \node[el,fill=warnpale,draw=warnred] at (0.25,-6.96)
    {\textbf{\textcolor{warnred}{Paid results}} --- bottom advertisements};
  \node[el,fill=white,draw=rulegrey] at (0.25,-7.74)
    {\textbf{Related searches}};

  \draw[rulegrey,line width=0.9pt,rounded corners=3pt] (7.3,-0.72) rectangle (10.6,-4.3);
  \node[anchor=north,align=center,text width=3.0cm,font=\sffamily\scriptsize,
        text=slate] at (8.95,-0.95)
    {\textbf{Knowledge panel}\\[3pt] entity information drawn from structured
     sources and the knowledge graph};

  \node[anchor=north west,align=left,text width=3.3cm,font=\sffamily\scriptsize,
        text=slatelight] at (7.3,-4.7)
    {\textbf{Three result types}\\[3pt]
     \textcolor{warnred}{paid}\\
     \textcolor{slate}{organic}\\
     \textcolor{deepblue}{universal / vertical}\\ {\scriptsize (images, video,
     news, local, shopping)}};
\end{tikzpicture}
\caption[Anatomy of a modern search engine results page]%
        {A modern SERP. Because paid listings, the featured snippet, the local
         pack and the knowledge panel can all precede the first ordinary organic
         result, ranking first organically no longer means appearing first on the
         page.}
\label{fig:serp}
\end{figure}

\noindent Main elements and result types on a modern SERP:

\begin{itemize}
\item \sfb{Paid results} --- search advertisements at the top and bottom,
      shopping and product listing advertisements, local service advertisements.
\item \sfb{Organic results} --- the ranked list of web pages, each with a title
      tag, a URL and a meta description\ix{meta description}.
\item \sfb{Featured snippet}, known as ``position zero'' --- a direct answer
      extracted from a page.
\item \sfb{Knowledge panel and knowledge graph} --- entity information drawn
      from structured sources.
\item \sfb{Local pack or map pack} --- three local businesses with map, ratings
      and directions.
\item \sfb{People Also Ask} --- expandable related questions, and an excellent
      content-research source.
\item \sfb{Rich results} --- reviews, star ratings, FAQs, recipes and events,
      driven by structured data markup.
\item \sfb{Image, video and news carousels}, and related searches at the foot of
      the page.
\end{itemize}

The three main types of result are therefore \textbf{paid}, \textbf{organic} and
\textbf{universal or vertical} results, the last covering images, video, news,
local and shopping.

\section{The Shift From Keywords to Intent}

Since 2015, Google's \keytermas{RankBrain}{RankBrain} and
\keytermas{Hummingbird}{Hummingbird} algorithms have focused on understanding
user intent rather than matching keyword strings. They expect high-quality
semantic data and well-structured outlines, not keyword stuffing.

\begin{keybox}[E-E-A-T]
Google evaluates content on \textbf{Experience, Expertise, Authoritativeness and
Trustworthiness}. Alongside this runs the link principle: \textbf{link gravity
beats quantity} --- a large quantity of \emph{quality} links is the goal, not
simply a large quantity of links.
\end{keybox}

\begin{warnbox}[The Post-Panda Reality]
Keyword stuffing is dead. Google's sole objective is surfacing the most relevant
answer to a query. If the query is ``how to make banana bread'', the content must
provide the best, most readable, fastest answer.

\smallskip
White-hat optimisation therefore means writing naturally, organising flow well,
and earning backlinks by creating remarkable, link-worthy data or insight.
\end{warnbox}

\section{The Six-Step SEO Process}

An SEO project follows six steps. Note that SEO does not begin and end with
initial work: ongoing success requires monitoring results and building meaningful
content continually.

\begin{figure}[htbp]
\centering
\begin{tikzpicture}[font=\sffamily\footnotesize,
    st/.style={draw=ocre,line width=0.85pt,fill=ocrepale,rounded corners=3pt,
               minimum height=1.5cm,text width=1.72cm,align=center,
               font=\sffamily\bfseries\scriptsize,text=slate,inner sep=3pt}]
  \node[st] (s1) at (0,0)     {1\\ Keyword research};
  \node[st] (s2) at (1.98,0)  {2\\ Reporting \& goal setting};
  \node[st] (s3) at (3.96,0)  {3\\ Content building};
  \node[st] (s4) at (5.94,0)  {4\\ Page optimization};
  \node[st] (s5) at (7.92,0)  {5\\ Social \& link building};
  \node[st] (s6) at (9.90,0)  {6\\ Follow-up reporting};
  \foreach \a/\b in {s1/s2,s2/s3,s3/s4,s4/s5,s5/s6}{
    \draw[-{Latex[length=1.9mm]},ocre,line width=0.85pt] (\a)--(\b);}
  \draw[-{Latex[length=2.1mm]},ocredark,line width=0.85pt,dashed]
    (s6.south) -- ++(0,-0.65) -| (s1.south);
  \node[anchor=north,text=ocredark,font=\sffamily\scriptsize\itshape]
    at (4.95,-1.5) {SEO is a cycle: results are monitored and content built continually};
\end{tikzpicture}
\caption[The six-step SEO process]%
        {The six-step process. Step~6 returns to step~1: an SEO project has no
         completion date, only measurement intervals.}
\label{fig:seoprocess}
\end{figure}

\begin{mmlongtable}{The six steps in detail}{tab:seosteps}%
{@{}P{0.22\textwidth}P{0.71\textwidth}@{}}
\rowcolor{ocre}\thd{Step} & \thd{Activity}\\
\midrule
\endfirsthead
\rowcolor{ocre}\thd{Step} & \thd{Activity}\\
\midrule
\endhead
\rlab{1. Keyword research} & Identify a group of keyword phrases balancing high
usage by searchers with relatively low competition. Also perform competitive
research --- a thorough analysis against the site's seven to ten biggest
competitors using metrics such as indexed content, inbound links, domain age and
social following, to gauge the starting position and identify priorities.\\
\rowcolor{tablealt}
\rlab{2. Reporting and goal setting} & Establish the site's current position in
the search engines and its traffic baseline --- which engines, which phrases,
bounce rates, most popular content. Set measurable goals tied to specific
business objectives.\\
\rlab{3. Content building} & Content is king. High-volume, high-quality content
gives users a reason to stay and return, and gives search engines more
information to store, which translates directly into rankings.\\
\rowcolor{tablealt}
\rlab{4. Page optimization} & On-page work: page titles, text-based navigation,
prominence of targeted keyword phrases, site map, ALT and META data, and cleaning
up the code.\\
\rlab{5. Social and link building} & Establish a social media presence to share
content --- fish where the fish are --- and build high-quality inbound links.
Google views links from high-quality sites as votes for your site.\\
\rowcolor{tablealt}
\rlab{6. Follow-up reporting and analysis} & Repeat the initial reporting at
regular intervals after optimisation. Compare rankings, traffic, social signals
and other key metrics against pre-optimisation levels to produce measurable
results.\\
\end{mmlongtable}

\section{SEO Goals, and White Hat Against Black Hat}

\begin{itemize}
\item The ultimate goal is to \sfb{increase the site's usability} so that the
      right people are brought in from the search engines. The purpose of the
      site must be clearly defined so that you can ensure the site achieves it.
\item When search visitors find that the site meets their expectations after
      clicking through, usability is solid. If they are disappointed, the site
      has missed its mark.
\item \sfb{Follow the best strategies for ranking}, which requires understanding
      how search engines want things done --- familiarise yourself with their
      terms and conditions and webmaster guidelines.
\item \sfb{Use white-hat techniques}, those that fall in line with what search
      engines define as best practice.
\end{itemize}

\begin{warnbox}[Black-Hat Techniques and Their Cost]
Keyword stuffing, cloaking, link farms, hidden text and doorway pages produce
short-term gains and long-term penalties. A manual action or algorithmic penalty
can remove a site from the index entirely, and recovery takes far longer than the
gain ever lasted.
\end{warnbox}

\begin{summarybox}
SEO improves visibility and ranking in organic results to raise the quantity and
quality of traffic; its objective is relevance as well as visibility. It matters
because organic search delivers about 53\% of all web traffic, carries the
highest buyer intent\ix{buyer intent}, compounds indefinitely and costs nothing per click. SEO is
organic, PPC is paid, SEM is the umbrella over both. A modern SERP carries paid
results, organic results and universal results, plus featured snippets,
knowledge panels, local packs, People Also Ask and rich results. Since 2015
RankBrain and Hummingbird have ranked on intent rather than string matching, and
content is judged on E-E-A-T. The process runs six steps --- keyword research,
reporting and goals, content building, page optimisation, social and link
building, follow-up analysis --- and cycles. White-hat technique follows engine
guidelines; black-hat technique trades short gains for long penalties.
\end{summarybox}

%%%%%%%%%%%%%%%%%%%%%%%%%%%%%%%%%%%%%%%%%%%%%%%%%%%%%%%%%%%%%%%%%%%%%%%%%%%%%%
\chapter{Keywords and the SEO Content Plan}
%%%%%%%%%%%%%%%%%%%%%%%%%%%%%%%%%%%%%%%%%%%%%%%%%%%%%%%%%%%%%%%%%%%%%%%%%%%%%%

\syllabusstrip{Keywords and SEO Content Plan.}

\begin{learningoutcomes}
\begin{itemize}
\item Why keyword research matters.
\item Keyword types by intent and by length.
\item The keyword sweet spot.
\item The ten-step SEO content plan, and keyword cannibalisation.
\end{itemize}
\end{learningoutcomes}

\section{Why Keyword Research Matters}

\begin{exambox}[Model Answer --- 2 Marks]
\pyq{CIE-II, 7 May 2025; SEE, Oct/Nov 2023 --- 2 M}
\keytermas{Keyword research}{keyword research} is the process of identifying the
phrases that real users type into search engines, together with their search
volume, competition and intent.

\smallskip
It matters because without it content is optimised for phrases nobody searches,
effort is spent on terms that cannot realistically be ranked for, and traffic
arrives with no intent to buy.
\end{exambox}

\section{Keyword Types}

\subsection{By intent}

\begin{mmlongtable}{Keyword types by search intent}{tab:keywordintent}%
{@{}P{0.24\textwidth}P{0.32\textwidth}P{0.37\textwidth}@{}}
\rowcolor{ocre}\thd{Intent type} & \thd{What the user wants} & \thd{Example}\\
\midrule
\endfirsthead
\rowcolor{ocre}\thd{Intent type} & \thd{What the user wants} & \thd{Example}\\
\midrule
\endhead
\rlab{Informational} & To learn something & ``what is content marketing''\\
\rowcolor{tablealt}
\rlab{Navigational} & To reach a specific site & ``RVCE login''\\
\rlab{Commercial investigation} & To compare before buying & ``best trekking
shoes under 5000''\\
\rowcolor{tablealt}
\rlab{Transactional} & To buy now & ``buy trekking shoes online''\\
\end{mmlongtable}

\subsection{By length}

\begin{figure}[htbp]
\centering
\begin{tikzpicture}[font=\sffamily\footnotesize]
  \fill[ocre!55] (0,0) rectangle (2.1,2.55);
  \node[text=white,font=\sffamily\bfseries\scriptsize,align=center] at (1.05,2.15) {HEAD};
  \node[text=slate,font=\sffamily\scriptsize,align=center,text width=1.95cm] at (1.05,1.15)
    {1--2 words\\ huge volume\\ brutal competition\\ vague intent};

  \fill[ocre!35] (2.3,0) rectangle (4.9,1.75);
  \node[text=slate,font=\sffamily\bfseries\scriptsize,align=center] at (3.6,1.45) {BODY};
  \node[text=slate,font=\sffamily\scriptsize,align=center,text width=2.4cm] at (3.6,0.78)
    {2--3 words\\ moderate volume\\ and competition};

  \fill[ocre!16] (5.1,0) rectangle (11.2,0.95);
  \node[text=slate,font=\sffamily\bfseries\scriptsize,align=center] at (8.15,0.74) {LONG TAIL};
  \node[text=slate,font=\sffamily\scriptsize,align=center,text width=5.9cm] at (8.15,0.3)
    {4+ words \ \textbullet\ \ low volume each but collectively the majority of
     all searches \ \textbullet\ \ low competition, sharp intent, high conversion};

  \draw[-{Latex[length=2.4mm]},slate,line width=0.9pt] (0,-0.4) -- (11.2,-0.4);
  \node[anchor=north,text=slatelight,font=\sffamily\scriptsize] at (5.6,-0.5)
    {increasing phrase length \ \textbullet\ \ decreasing competition \ \textbullet\ \ increasing intent};
\end{tikzpicture}
\caption[Keywords by length: the search demand curve]%
        {The search demand curve. The long tail holds little volume per phrase
         and most of the volume in total, which is why a site with low authority
         competes there rather than at the head.}
\label{fig:longtail}
\end{figure}

\section{Finding the Keyword Sweet Spot}

A keyword is worth targeting only where three conditions overlap.

\begin{figure}[htbp]
\centering
\begin{tikzpicture}[font=\sffamily\footnotesize]
  \begin{scope}[opacity=0.40,transparency group]
    \fill[ocre]     (-1.3,0.6) circle (2.25);
    \fill[deepblue] ( 1.3,0.6) circle (2.25);
    \fill[greenok]  ( 0.0,-1.55) circle (2.25);
  \end{scope}
  \node[text=ocredark,font=\sffamily\bfseries\scriptsize,align=center]
    at (-2.55,1.75) {HIGH SEARCH\\ VOLUME};
  \node[text=deepblue,font=\sffamily\bfseries\scriptsize,align=center]
    at ( 2.55,1.75) {LOW SEO\\ DIFFICULTY};
  \node[text=greenok,font=\sffamily\bfseries\scriptsize,align=center]
    at ( 0.0,-3.4)  {COMMERCIAL INTENT};
  \node[text=slate,font=\sffamily\scriptsize,align=center,text width=2.1cm]
    at (-2.25,0.5) {roughly 250\\ to 1,000+\\ monthly};
  \node[text=slate,font=\sffamily\scriptsize,align=center,text width=2.1cm]
    at ( 2.25,0.5) {under about\\ 40 KD/SD};
  \node[text=slate,font=\sffamily\scriptsize,align=center,text width=2.4cm]
    at ( 0.0,-2.1) {high CPC, buyer-\\ specific phrasing};
  \node[text=white,font=\sffamily\bfseries\scriptsize,align=center]
    at (0,-0.32) {THE SWEET\\ SPOT};
\end{tikzpicture}
\caption[The keyword sweet spot]%
        {Target only the central overlap. Two out of three is a trap: volume
         without intent is vanity, and intent without achievability is wasted
         effort.}
\label{fig:sweetspot2}
\end{figure}

\begin{itemize}
\item \sfb{High search volume} --- a minimum of roughly 250 to 1,000 or more
      monthly searches.
\item \sfb{Low SEO difficulty} --- under about a 40 difficulty score, so that
      ranking is realistically achievable.
\item \sfb{Commercial intent} --- a high cost-per-click and buyer-specific
      phrasing, indicating the searcher intends to transact.
\end{itemize}

\begin{warnbox}[Traffic Without Intent Is Useless]
Do not optimise for clicks; optimise for revenue.
\end{warnbox}

\section{Building an SEO Content Plan}

\begin{enumerate}
\item \sfb{Define business objectives first} --- what does the site actually need
      (sales, leads, sign-ups)?
\item \sfb{Research keywords} and cluster them into topics rather than isolated
      phrases.
\item \sfb{Map intent to funnel stage} --- informational keywords become blog
      content; commercial and transactional keywords become product and
      comparison pages.
\item \sfb{Audit existing content} --- decide what to keep, update, consolidate
      or retire.
\item \sfb{Identify content gaps} against competitors --- topics with demand
      where nobody ranks well.
\item \sfb{Design the site architecture} --- a pillar page per topic cluster,
      with supporting articles internally linked to it.
\item \sfb{Assign one primary keyword per page.}
\item \sfb{Build the editorial calendar} with owners, deadlines and target
      keywords.
\item \sfb{Set measurable targets} --- rankings, impressions, organic sessions,
      assisted conversions.
\item \sfb{Review and iterate} monthly.
\end{enumerate}

\begin{warnbox}[Keyword Cannibalisation]
Never let two of your own pages compete for the same term.
\keytermas{Keyword cannibalisation}{keyword cannibalisation} splits link equity
and click-through between two weaker pages instead of concentrating it in one
strong page, and it leaves the search engine to choose which of your pages to
rank --- a choice it frequently gets wrong. One primary keyword per page,
enforced in the plan, prevents it.
\end{warnbox}

\begin{summarybox}
Keyword research identifies the phrases users actually type, with their volume,
competition and intent; without it, effort goes to phrases nobody searches or
nobody buys from. Intent divides into informational, navigational, commercial
investigation and transactional; length divides into head, body and long tail,
with the long tail holding most total volume at low competition and high
conversion. Target only where volume, achievable difficulty and commercial intent\ix{commercial intent}
overlap. The content plan runs ten steps from business objectives through topic
clustering, intent mapping, content audit, gap analysis, pillar architecture, one
keyword per page, calendar and targets to monthly review.
\end{summarybox}

%%%%%%%%%%%%%%%%%%%%%%%%%%%%%%%%%%%%%%%%%%%%%%%%%%%%%%%%%%%%%%%%%%%%%%%%%%%%%%
\chapter{SEO and Business Objectives}
%%%%%%%%%%%%%%%%%%%%%%%%%%%%%%%%%%%%%%%%%%%%%%%%%%%%%%%%%%%%%%%%%%%%%%%%%%%%%%

\syllabusstrip{SEO \& Business Objectives.}

\begin{learningoutcomes}
\begin{itemize}
\item How to translate each business objective into an SEO focus and a KPI.
\item Why SEO must be designed into website development rather than retro-fitted.
\end{itemize}
\end{learningoutcomes}

\section{Aligning SEO With Business Objectives}

\begin{mmlongtable}{SEO focus and KPI by business objective}{tab:seoobjectives}%
{@{}P{0.22\textwidth}P{0.42\textwidth}P{0.29\textwidth}@{}}
\rowcolor{ocre}\thd{Business objective} & \thd{SEO focus} & \thd{Primary KPI}\\
\midrule
\endfirsthead
\rowcolor{ocre}\thd{Business objective} & \thd{SEO focus} & \thd{Primary KPI}\\
\midrule
\endhead
\rlab{Increase online sales} & Transactional and product keywords; optimised
category and product pages; rich results with reviews and price & Organic
revenue; organic conversion rate\\
\rowcolor{tablealt}
\rlab{Generate leads} & Commercial-investigation keywords; comparison and ``best
X for Y'' content; gated assets & Organic leads; cost per lead\\
\rlab{Build brand authority} & Informational pillar content; E-E-A-T signals;
digital PR and backlinks & Impressions; non-branded keyword coverage; referring
domains\\
\rowcolor{tablealt}
\rlab{Reduce paid advertising dependence} & Rank organically for the terms
currently being bought & Share of traffic from organic against paid\\
\rlab{Enter a new market} & Local and language-specific SEO; localised content
and listings & Local pack visibility; regional organic sessions\\
\end{mmlongtable}

\section{Integrating SEO Into Website Development}

\begin{exambox}[Model Answer --- 2 Marks]
``Why is it important to integrate SEO in website development?''
\pyq{SEE, Jun/Jul 2025 --- 2 M}

\begin{itemize}
\item Site architecture, URL structure and navigation are extremely expensive to
      change after launch but nearly free to get right during the build.
\item Rendering method, page speed and Core Web Vitals\ix{Core Web Vitals} are determined by
      development decisions, not by marketing decisions.
\item Retro-fitting SEO usually forces URL changes, which risk losing
      accumulated ranking authority.
\item Content structure --- heading hierarchy, schema, internal linking --- must
      be designed into the templates.
\item Tracking and analytics must be instrumented before launch, not after
      traffic has been lost.
\item It avoids launching with pages blocked by \texttt{robots.txt} or left
      \texttt{noindex} from staging --- a classic and costly error.
\end{itemize}
\end{exambox}

\begin{summarybox}
Every business objective implies a different SEO focus and a different KPI:
sales imply transactional keywords and organic revenue; leads imply comparison
content and cost per lead\ix{cost per lead}; authority implies pillar content and referring
domains; reducing paid dependence implies ranking for bought terms; new-market
entry implies local and language SEO. SEO belongs in the build, not after it,
because architecture, URLs, rendering, speed, template structure and tracking are
all fixed at development time, and retro-fitting risks the accumulated authority
of existing URLs.
\end{summarybox}

%%%%%%%%%%%%%%%%%%%%%%%%%%%%%%%%%%%%%%%%%%%%%%%%%%%%%%%%%%%%%%%%%%%%%%%%%%%%%%
\chapter{Writing SEO Content}
%%%%%%%%%%%%%%%%%%%%%%%%%%%%%%%%%%%%%%%%%%%%%%%%%%%%%%%%%%%%%%%%%%%%%%%%%%%%%%

\syllabusstrip{Writing SEO Content.}

\begin{learningoutcomes}
\begin{itemize}
\item The anatomy of a top-ranking page.
\item The nine-point method for integrating keywords naturally.
\item Seven considerations when writing SEO-friendly content.
\end{itemize}
\end{learningoutcomes}

\section{Anatomy of a Top-Ranking Page}

\begin{itemize}
\item \sfb{H1 title} --- click-worthy, with the exact-match keyword.
\item \sfb{Introduction} --- comprehensive, answering the user's core intent
      immediately rather than after four hundred words of preamble.
\item \sfb{Structure} --- H2 and H3 subheadings containing semantically related
      long-tail keyword\ix{long-tail keyword} variants.
\item \sfb{Formatting} --- short, scannable paragraphs and bulleted lists.
\item \sfb{Rich media} --- embedded videos and infographics to increase dwell
      time.
\item \sfb{Architecture} --- internal links to keep users inside your ecosystem,
      plus credible outbound links to build trust.
\end{itemize}

\section{Integrating Keywords Naturally}

\begin{exambox}[The Scenario Question]
``You are writing a blog post about a complex technical topic. Describe how you
would integrate relevant keywords naturally into the content while ensuring it
remains engaging and informative for readers.''
\pyq{CIE-II, 24 Jul 2024 --- 10 M}
\end{exambox}

\begin{enumerate}
\item \sfb{Write for the reader first, then optimise.} Draft the piece to
      genuinely answer the question, then place keywords during editing.
\item \sfb{Place the primary keyword in the high-value positions} --- the title
      tag, the H1, the first hundred words, one H2, the URL slug, the meta
      description and one image ALT attribute.
\item \sfb{Use semantic variants and synonyms} rather than repeating the exact
      phrase; modern algorithms understand meaning, not string matching.
\item \sfb{Respect keyword density.} Industry practice is roughly 3--5\% for the
      focus phrase --- enough to signal the topic, not enough to read as spam.
\item \sfb{Answer the sub-questions surfaced in ``People Also Ask''}; each
      becomes an H2 or H3 that naturally carries a long-tail variant.
\item \sfb{Explain jargon on first use}, and use analogies, examples, diagrams
      and code snippets to keep a technical topic engaging.
\item \sfb{Break up the text} with subheadings, lists, tables and visuals every
      few paragraphs.
\item \sfb{Link internally} with keyword-rich but descriptive anchor text\ix{anchor text}.
\item \sfb{Read it aloud.} If a sentence sounds unnatural because of a keyword,
      rewrite it --- user signals now outweigh keyword placement.
\end{enumerate}

\section{Key Considerations}

\begin{exambox}[Two Considerations --- 2 Marks]
\pyq{CIE-II Quiz, 24 Jul 2024 --- 2 M}
(1) \textbf{Match the search intent} --- the format and depth must answer what
the user actually wants. (2) \textbf{Integrate keywords naturally} in the title,
H1, opening paragraph and subheadings using semantic variants, at roughly 3--5\%
density, without stuffing. Also acceptable: a unique title tag and meta
description; a scannable structure.
\end{exambox}

\begin{itemize}
\item \sfb{Search intent match} --- the format must match what already ranks: a
      list, a guide, a comparison, a tool.
\item \sfb{Originality and depth} --- content must be \emph{better} than what
      currently occupies the SERP, not merely equivalent.
\item \sfb{Readability} --- short sentences, plain language, scannable structure.
\item \sfb{Freshness} --- update statistics and examples; stale content loses
      rankings.
\item \sfb{Unique title tags and meta descriptions} on every page.
\item \sfb{Accessibility} --- descriptive ALT text on every meaningful image.
\item \sfb{Comprehensiveness} --- cover the topic fully, so the user has no
      reason to return to the SERP.
\end{itemize}

\begin{summarybox}
A top-ranking page has a keyword-bearing H1, an introduction that answers the
intent immediately, H2 and H3 structure carrying semantic variants, scannable
formatting, rich media for dwell time\ix{dwell time}, and both internal and credible outbound
links. Integrate keywords by writing for the reader first and optimising second,
placing the primary phrase in the high-value positions, using semantic variants
at roughly 3--5\% density, answering People Also Ask questions as subheadings,
explaining jargon, breaking up text, linking internally with descriptive anchors,
and reading aloud to catch unnatural phrasing. Content must match intent, exceed
what already ranks, read easily, stay fresh, carry unique metadata, be accessible
and be complete.
\end{summarybox}

%%%%%%%%%%%%%%%%%%%%%%%%%%%%%%%%%%%%%%%%%%%%%%%%%%%%%%%%%%%%%%%%%%%%%%%%%%%%%%
\chapter{On-Site and Off-Site SEO}
%%%%%%%%%%%%%%%%%%%%%%%%%%%%%%%%%%%%%%%%%%%%%%%%%%%%%%%%%%%%%%%%%%%%%%%%%%%%%%

\syllabusstrip{On-site \& off-site SEO.}

\begin{learningoutcomes}
\begin{itemize}
\item On-page SEO and its five key elements.
\item Nine interlinking best practices.
\item Technical SEO --- the defensive foundation.
\item Off-page SEO, backlinks and brand signals.
\item The on-site against off-site comparison.
\item The House of SEO --- how the three layers combine.
\end{itemize}
\end{learningoutcomes}

\section{On-Site (On-Page) SEO}

\begin{definitionbox}[On-Page SEO]
\keytermas{On-page SEO}{on-page SEO} refers to how well a website's own content
and code are presented to search engines. It involves ensuring that a particular
webpage is structured so that it is found for given keywords and key phrases. It
improves search ranking and the overall readability of the site, and --- unlike
off-page work --- it can be improved immediately by fixing incorrect elements on
the page.
\end{definitionbox}

\begin{mmlongtable}{The five key elements of on-page SEO}{tab:onpage}%
{@{}P{0.16\textwidth}P{0.77\textwidth}@{}}
\rowcolor{ocre}\thd{Element} & \thd{What to do}\\
\midrule
\endfirsthead
\rowcolor{ocre}\thd{Element} & \thd{What to do}\\
\midrule
\endhead
\rlab{1. Page copy} & Produce original, unique, high-quality, relevant content
continuously. Focus each piece on a single primary keyword or key phrase,
maintaining a keyword density of about 3--5\%, mixing primary and secondary
phrases. Quality over quantity --- depth and content intensity have outperformed
sheer length since the 2011 algorithm update.\\
\rowcolor{tablealt}
\rlab{2. Title tags} & Arguably the most important element of the ``big three'',
the others being page copy and inbound links. The title tag supplies the
clickable link in the search result. Google limits page titles to about 70
characters, so titles must be keyword-relevant yet concise. Order: primary
keyword, then secondary keyword, then branded keyword at the end --- except on
the home page.\\
\rlab{3. Meta data} & A well-written description summarising the page, limited to
155--160 characters. Meta keywords lost ranking influence after Google's
September 2009 declaration because of spam abuse, though crawlers still read them
for topicality. The meta description \emph{cannot} influence rankings but
strongly influences click-through rate --- it works as advertising copy for the
organic result, and Google highlights matching keywords in bold.\\
\rowcolor{tablealt}
\rlab{4. Heading tags} & Define the content section by section, like traditional
headings and subheadings. There should be one \texttt{<h1>} per page containing
the most relevant key phrase. Tags run to \texttt{<h6>}, though general practice
uses up to \texttt{<h3>}. Place other important key phrases in \texttt{<h2>} and
\texttt{<h3>}.\\
\rlab{5. Interlinking} & Linking one page of your site to related pages provides
context to both search engines and readers.\\
\end{mmlongtable}

\subsection{Interlinking best practices}

\begin{enumerate}
\item Include links in the main content of each page.
\item Paragraph links carry the most weight.
\item Use keyword-rich anchor text.
\item Avoid non-descriptive anchor text such as ``read more'' or ``click here''.
\item Link to relevant, deep pages --- not only the home page.
\item Use breadcrumb navigation on every page.
\item Monitor inbound links through Google Search Console\ix{Google Search Console}.
\item Avoid multiple links to the same page from a single page.
\item Fewer links means more authority per link.
\end{enumerate}

\section{Technical SEO --- The Defence}

\begin{quotebox}
SEO is like a sports team. To win, both defence and offence are required.
\end{quotebox}

Technical SEO is the foundation of the house: the architecture that allows the
site to be crawled and indexed at all.

\begin{itemize}
\item Site speed and \sfb{Core Web Vitals} --- LCP, INP and CLS.
\item Mobile friendliness and responsive design\ix{responsive design}.
\item HTTPS security.
\item An XML sitemap and a correct \texttt{robots.txt}.
\item Crawlability and indexing\ix{indexing} --- no orphan pages, no broken links, no
      redirect chains.
\item Clean URL structure --- short, readable, keyword-bearing.
\item Canonical tags, to prevent duplicate-content dilution.
\item Structured data and schema markup, for rich results.
\item Avoiding intrusive interstitials and pop-ups that block content on mobile.
\end{itemize}

\begin{exambox}[Two Technical Aspects --- 2 Marks]
``Name two technical aspects of a website that can significantly impact its
organic search ranking.''
\pyq{CIE-II Quiz, 24 Jul 2024 --- 2 M}
Best two answers: \textbf{site loading speed} (Core Web Vitals) and
\textbf{mobile responsiveness}. Also acceptable: HTTPS security, crawlability and
indexability, XML sitemap, URL structure.
\end{exambox}

\section{Off-Site (Off-Page) SEO}

\begin{definitionbox}[Off-Page SEO]
\keytermas{Off-page SEO}{off-page SEO} refers to a website's overall authority on
the web, determined by what other websites say about it. It is a long-term
process. Put simply, off-page is about your online reputation. Its central
activity is acquiring \keytermas{backlinks}{backlink} from authority sites in
your niche --- backlinks are the currency of any off-page strategy. Unlike
on-page work, off-page effort is not visible on the webpage itself; it does the
background work for a better search result.
\end{definitionbox}

\subsection{The two halves of off-page work}

\noindent\sfb{1. Acquiring backlinks.} Search engines treat link popularity as a
key ranking factor, though it is no longer the top factor, because it can be
manipulated. After the Google Panda and Penguin updates, many effective
old-school practices became obsolete and negatively affected large, high-ranking
websites. Engines now focus more on content quality and engagement level than on
the raw number of links. The success factor is not building a long list of
inbound links but building a trail of quality links.

\begin{keybox}[The Link Rule]
Getting links from multiple \emph{relevant} domains in your industry is the key.
\textbf{Link gravity beats quantity} --- a large quantity of quality links is the
goal.
\end{keybox}

\noindent\sfb{2. Building brand signals} --- brand mentions, reviews, ratings,
social sharing, PR coverage and demand generation, all of which contribute to
E-E-A-T.

\subsection{Off-page techniques}

\begin{itemize}
\item \sfb{Digital PR and link earning} --- publish original research, data or
      tools that journalists and bloggers cite naturally.
\item \sfb{Guest posting} on relevant, authoritative sites in the niche.
\item \sfb{Broken-link building} --- find dead links on relevant pages and offer
      your resource as a replacement.
\item \sfb{Business listings and citations} --- Google Business Profile, industry
      directories, consistent name, address and phone details.
\item \sfb{Online reviews and ratings} on Google, industry platforms and
      marketplaces.
\item \sfb{Social media sharing} --- not a direct ranking factor, but it produces
      the exposure from which links follow.
\item \sfb{Influencer and expert collaborations} --- interviews, quotes and
      co-created content.
\item \sfb{Forum and community participation} --- contributing genuine value in
      niche communities.
\item \sfb{Unlinked-mention reclamation} --- find places where the brand is
      mentioned without a link, and request one.
\end{itemize}

\section{On-Site Against Off-Site SEO}

\begin{keybox}[High Yield]
Asked both as a two-mark distinction and as a ten-mark analysis in the same
paper.
\pyq{CIE-II, 7 May 2025 --- 2 M and 10 M}
\end{keybox}

\begin{mmlongtable}{On-site against off-site SEO}{tab:onvsoff}%
{@{}P{0.15\textwidth}P{0.39\textwidth}P{0.39\textwidth}@{}}
\rowcolor{ocre}\thd{Basis} & \thd{On-Site (On-Page) SEO} & \thd{Off-Site (Off-Page) SEO}\\
\midrule
\endfirsthead
\rowcolor{ocre}\thd{Basis} & \thd{On-Site (On-Page) SEO} & \thd{Off-Site (Off-Page) SEO}\\
\midrule
\endhead
\rlab{Location of work} & Within your own website & Outside your website\\
\rowcolor{tablealt}
\rlab{What it controls} & How content and code are presented to search engines &
The site's authority and reputation across the web\\
\rlab{Primary target} & Humans and crawlers reading the page & Third-party
validators --- other sites and users\\
\rowcolor{tablealt}
\rlab{Core challenge} & Quality and formatting & Brand authority\\
\rlab{Degree of control} & Complete --- you own it & Limited --- others decide\\
\rowcolor{tablealt}
\rlab{Speed of effect} & Immediate; can be fixed today & Slow; a long-term
process\\
\rlab{Techniques} & Title tags, meta descriptions, heading tags, keyword
placement, content quality, internal linking, image ALT text, URL structure,
site speed, mobile friendliness & Backlink building, guest posting, digital PR,
social sharing, influencer collaboration, online reviews, business listings,
forum participation\\
\rowcolor{tablealt}
\rlab{Example} & Rewriting a product page title to ``Trekking Shoes for Men ---
Waterproof --- BrandName'' & Earning a link from a national outdoors magazine's
``best trekking gear'' article\\
\end{mmlongtable}

\section{The House of SEO}

\begin{figure}[htbp]
\centering
% Each line of each storey is a separately anchored node at a fixed y, so that
% text can never grow into the storey below it.
\begin{tikzpicture}[font=\sffamily\footnotesize,
    ln/.style={anchor=north,align=center,font=\sffamily\scriptsize,text=slate}]

  % ---- roof: off-page SEO
  \fill[greenpale]              (0,4.3) -- (5.6,6.1) -- (11.2,4.3) -- cycle;
  \draw[greenok,line width=1pt] (0,4.3) -- (5.6,6.1) -- (11.2,4.3) -- cycle;
  \node[ln,text=greenok,font=\sffamily\bfseries\small] at (5.6,5.72) {OFF-PAGE SEO};
  \node[ln,font=\sffamily\scriptsize\itshape]           at (5.6,5.26) {the roof --- the offence};
  \node[ln,text width=8.0cm]                            at (5.6,4.90)
    {link building, PR, ratings, demand generation, E-E-A-T};

  % ---- living floor: on-page SEO
  \fill[ocrepale]            (0.55,1.35) rectangle (10.65,4.3);
  \draw[ocre,line width=1pt] (0.55,1.35) rectangle (10.65,4.3);
  \node[ln,text=ocredark,font=\sffamily\bfseries\small] at (5.6,4.18) {ON-PAGE SEO};
  \node[ln,font=\sffamily\scriptsize\itshape]           at (5.6,3.74) {the living floor};
  \node[ln,text width=9.4cm]                            at (5.6,3.38)
    {relevant topics, title tags, H1--H6 headings, image ALT text, multimedia};
  \node[ln,text width=9.4cm]                            at (5.6,2.72)
    {\textbf{goal:} publish original, error-free content better than the SERP
     competitors};

  % ---- foundation: technical SEO
  \fill[deepbluepale]            (0.55,-1.15) rectangle (10.65,1.35);
  \draw[deepblue,line width=1pt] (0.55,-1.15) rectangle (10.65,1.35);
  \node[ln,text=deepblue,font=\sffamily\bfseries\small] at (5.6,1.23) {TECHNICAL SEO};
  \node[ln,font=\sffamily\scriptsize\itshape]           at (5.6,0.79) {the foundation --- the defence};
  \node[ln,text width=9.4cm]                            at (5.6,0.43)
    {URL structure, HTTPS, mobile friendliness, Core Web Vitals, fast hosting};
  \node[ln,text width=9.4cm]                            at (5.6,-0.23)
    {\textbf{goal:} complete crawlability, indexing and seamless UX
     infrastructure};

  % ---- reading of the diagram
  \node[ln,text=slatelight,font=\sffamily\scriptsize\itshape,text width=11cm]
    at (5.6,-1.45)
    {technical makes the site discoverable, on-page makes it relevant, off-page
     makes it authoritative --- and rankings are a function of all three
     together};
\end{tikzpicture}
\caption[The House of SEO]%
        {The House of SEO. The layers are not alternatives: authority cannot
         compensate for a site that cannot be crawled, and crawlability earns
         nothing without content worth ranking.}
\label{fig:houseofseo}
\end{figure}

\begin{mmlongtable}{The three layers compared}{tab:threelayers}%
{@{}P{0.147\textwidth}P{0.255\textwidth}P{0.255\textwidth}P{0.255\textwidth}@{}}
\rowcolor{ocre}\thd{} & \thd{Technical SEO} & \thd{On-Page SEO} & \thd{Off-Page SEO}\\
\midrule
\endfirsthead
\rowcolor{ocre}\thd{} & \thd{Technical SEO} & \thd{On-Page SEO} & \thd{Off-Page SEO}\\
\midrule
\endhead
\rlab{Position} & The foundation & The living floor & The roof\\
\rowcolor{tablealt}
\rlab{Primary target} & Search engine bots & Humans and algorithms &
Third-party validators\\
\rlab{Core challenge} & Discoverability and security & Quality and formatting &
Brand authority\\
\rowcolor{tablealt}
\rlab{The goal} & Complete crawlability, indexing and seamless UX
infrastructure & Publish well-written, original, error-free content better than
SERP competitors & Enhance brand awareness and recognition; acquire high-quality
backlinks\\
\rlab{Elements} & URL structure, HTTPS, mobile friendliness, Core Web Vitals,
fast hosting, avoiding intrusive interstitials & Relevant topics, title tags,
H1--H6 headings, image ALT text, multimedia & Link building, PR, ratings, demand
generation, E-E-A-T\\
\end{mmlongtable}

\begin{summarybox}
On-page SEO governs how a site's own content and code are presented, through five
elements --- page copy at 3--5\% keyword density, title tags under about 70
characters, meta descriptions of 155--160 characters that drive click-through
rather than ranking, one H1 per page with H2 and H3 beneath, and interlinking with
descriptive anchor text. Technical SEO is the defensive foundation: speed, Core
Web Vitals, mobile friendliness, HTTPS, sitemap, crawlability, clean URLs,
canonicals and schema. Off-page SEO is authority earned elsewhere, principally
through quality backlinks --- link gravity beats quantity --- plus brand signals.
On-page is controllable and immediate; off-page is uncontrollable and slow. The
House of SEO combines all three: technical makes a site discoverable, on-page
makes it relevant, off-page makes it authoritative.
\end{summarybox}

%%%%%%%%%%%%%%%%%%%%%%%%%%%%%%%%%%%%%%%%%%%%%%%%%%%%%%%%%%%%%%%%%%%%%%%%%%%%%%
\chapter{Optimizing Organic Search Ranking}
%%%%%%%%%%%%%%%%%%%%%%%%%%%%%%%%%%%%%%%%%%%%%%%%%%%%%%%%%%%%%%%%%%%%%%%%%%%%%%

\syllabusstrip{Optimize Organic Search Ranking.}

\begin{learningoutcomes}
\begin{itemize}
\item The purpose of optimising organic rankings.
\item Eleven metrics by which organic ranking is measured.
\item The nine-step SEO audit, and a worked diagnostic example.
\item Why good design is now good SEO.
\item The SEO toolset.
\end{itemize}
\end{learningoutcomes}

\section{Purpose}

\begin{exambox}[Model Answer --- 2 Marks]
\pyq{CIE-II, 7 May 2025 --- 2 M}
The purpose of optimising organic search rankings is to appear higher in unpaid
search results for queries with real commercial intent, thereby capturing
sustainable, high-quality, near-zero-marginal-cost traffic; reducing dependence
on paid advertising; and building the credibility that comes with ranking
organically rather than paying for position.
\end{exambox}

\section{Key Metrics}

\begin{mmlongtable}{Metrics for organic search performance}{tab:organicmetrics}%
{@{}P{0.30\textwidth}P{0.63\textwidth}@{}}
\rowcolor{ocre}\thd{Metric} & \thd{What it measures}\\
\midrule
\endfirsthead
\rowcolor{ocre}\thd{Metric} & \thd{What it measures}\\
\midrule
\endhead
\rlab{Average position and keyword rankings} & Where the site sits for its
target keywords\\
\rowcolor{tablealt}
\rlab{Organic impressions} & How often pages are shown in search results\\
\rlab{Organic clicks and organic CTR} & How often people actually click --- and
how compelling the title and meta description are\\
\rowcolor{tablealt}
\rlab{Organic sessions} & Volume of visitors arriving from unpaid search\\
\rlab{Keyword coverage} & The number of distinct keywords ranking in the top 3,
top 10 and top 100\\
\rowcolor{tablealt}
\rlab{Indexed pages} & How much of the site the engine has actually stored\\
\rlab{Referring domains and backlink quality} & Off-page authority growth\\
\rowcolor{tablealt}
\rlab{Domain and page authority} & Third-party authority scoring\\
\rlab{Bounce rate and dwell time} & Whether the traffic finds what it expected\\
\rowcolor{tablealt}
\rlab{Organic and assisted conversions} & Whether the traffic produces business
value\\
\rlab{Core Web Vitals} & Technical experience signals feeding into ranking\\
\end{mmlongtable}

\section{Conducting an SEO Audit}

\begin{keybox}[High Yield]
Examined together with the metrics of \S\,11.2.
\pyq{CIE-II, 24 Jul 2024 --- 10 M}
\end{keybox}

\begin{enumerate}
\item \sfb{Crawl the site} --- with Screaming Frog\ix{Screaming Frog}, SEMrush Site Audit or Ahrefs
      Site Audit --- to find broken links, redirect chains, duplicate titles and
      descriptions, missing H1s, thin pages and orphan pages.
\item \sfb{Check indexation} in Google Search Console. Are the important pages
      actually indexed? Are unimportant ones consuming crawl\ix{crawl} budget?
\item \sfb{Audit technical health} --- speed, Core Web Vitals, mobile usability,
      HTTPS, sitemap, \texttt{robots.txt}, canonical tags, structured data.
\item \sfb{Audit on-page elements} --- title tags, meta descriptions, heading
      hierarchy, keyword targeting, keyword cannibalisation, image ALT text,
      internal linking.
\item \sfb{Audit content} --- relevance, depth, freshness, originality, intent
      match, and gaps against competitors.
\item \sfb{Audit backlinks} --- referring domains, anchor-text distribution,
      toxic links to disavow, and the link gap against competitors.
\item \sfb{Benchmark competitors} on the same metrics, to establish relative
      position.
\item \sfb{Prioritise findings by impact against effort}, and produce a ranked
      action plan.
\item \sfb{Implement, then re-measure} at defined intervals against the
      pre-optimisation baseline.
\end{enumerate}

\begin{examplebox}[Worked Diagnostic: an Institution That Cannot Rank for ``Online Courses'']
\pyq{CIE-II, 24 Jul 2024 --- 10 M}

\textbf{Likely issues.} Targeting head keywords far above the site's authority;
thin course pages with duplicate descriptions; no topic-cluster architecture;
slow, unoptimised pages; no schema markup for courses; very few referring domains
compared with established competitors; and no student-facing informational
content capturing early-funnel search.

\textbf{Recommended actions.}
\begin{enumerate}
\item Shift targeting to achievable long-tail phrases --- ``online data
      analytics course for working professionals in Bengaluru''.
\item Build a pillar page per subject area with supporting articles internally
      linked.
\item Rewrite every course page with unique titles, meta descriptions, outcomes,
      fees, faculty and FAQs.
\item Add Course and FAQ schema for rich results.
\item Fix speed and mobile usability.
\item Earn links through original placement-outcome research, faculty guest
      articles and local news coverage.
\item Build reviews on Google Business Profile and education directories.
\item Measure keyword coverage, organic sessions and enquiry conversions
      monthly.
\end{enumerate}
\end{examplebox}

\section{The Human--Algorithm Symbiosis}

\begin{keybox}[Why Good Design Is Now Good SEO]
Algorithms care about crawlability, keywords, metadata and links. Humans care
about scanning, speed, usability and satisfaction. The \textbf{overlap} is Core
Web Vitals --- speed and stability --- and E-E-A-T --- trust and authority.

\smallskip
Search engines have evolved to mimic human preferences. Good web usability is now
a fundamental requirement for good SEO, not a separate concern.
\end{keybox}

\section{SEO Tools}

\begin{mmlongtable}{SEO tools by category}{tab:seotools}%
{@{}P{0.24\textwidth}P{0.69\textwidth}@{}}
\rowcolor{ocre}\thd{Category} & \thd{Tools}\\
\midrule
\endfirsthead
\rowcolor{ocre}\thd{Category} & \thd{Tools}\\
\midrule
\endhead
\rlab{All-in-one suites} & SEMrush, Ahrefs, Moz Pro, Ubersuggest, SE Ranking\\
\rowcolor{tablealt}
\rlab{Google's own tools} & Google Search Console, Google Analytics 4, Keyword
Planner, PageSpeed Insights, Trends, Rich Results Test, Tag Manager\\
\rlab{Keyword research} & Keyword Planner, SEMrush Keyword Magic Tool, Ahrefs
Keywords Explorer, AnswerThePublic, AlsoAsked\\
\rowcolor{tablealt}
\rlab{Technical crawling} & Screaming Frog SEO Spider, Sitebulb, DeepCrawl\\
\rlab{Backlink analysis} & Ahrefs, Majestic, Moz Link Explorer, SEMrush Backlink
Audit\\
\rowcolor{tablealt}
\rlab{Rank tracking} & SEMrush Position Tracking, AccuRanker, SERPWatcher\\
\rlab{Content optimisation} & Surfer SEO, Clearscope, Yoast SEO, Rank Math\\
\rowcolor{tablealt}
\rlab{Competitive intelligence} & SimilarWeb, SpyFu, BuzzSumo\\
\end{mmlongtable}

\begin{summarybox}
Optimising organic rankings captures sustainable high-intent traffic at near-zero
marginal cost, reduces paid dependence and confers the credibility of an organic
listing. Performance is measured on rankings, impressions, clicks and CTR,
sessions, keyword coverage, indexed pages, referring domains, authority scores,
bounce\ix{bounce} and dwell, conversions and Core Web Vitals. The audit runs nine steps from
crawl and indexation through technical, on-page, content and backlink review to
competitor benchmarking, impact-versus-effort prioritisation, implementation and
re-measurement. Because search engines have evolved to mimic human preference,
good usability is now a requirement of good SEO rather than a separate concern.
\end{summarybox}
